% Notes of the Montreal seminar, 1970 (cote 129, typed part): archive cote 129, site page 53 (PDF page 54).
% typescript. Draft compared: none.
% AI-transcribed by Claude Opus 5.5 (claude-opus-5-5, adaptive thinking, effort high), Message Batches API, batch mtl129_full (msgbatch_01AGnDKHwy6ATCM7ehUa6263), 2026-11-13; not yet checked by eye.
\scanpage{129/53}
\mspage{\hand{II.25}}

6.4. \uline{Le foncteur} $\mathbb{E}$ \uline{est adjoint à gauche de} D*

On veut construire un isomorphisme fonctoriel en G et M
\[ \mathrm{Hom}_{k\text{-}Gr}(G,\mathbb{E}(M)) \overset{\sim}{\to} \mathrm{Hom}_{\hand{$\mathcal{D}$}}(M,\ D^*(G)) \ . \]
Si l'on pose
\margin{XX}
\[ \mathrm{Prim}(\deleted{(}G) = \mathrm{Ker}\ \{\underset{\rightarrow}{\mathbb{W}}(G) \xrightarrow{\ \Delta^*-pr_1^*-pr_2^*\ } \underset{\rightarrow}{\mathbb{W}}(G \times G)\}\deleted{)}\ , \]
on vérifie que l'on a :
\[ D^*(G) = \mathrm{Hom}_{k\text{-}Gr}(G,\underset{\rightarrow}{\mathbb{W}}) \overset{\sim}{\to} \mathrm{Prim}(G)\ . \]
Par ailleurs, d'après la définition même de $\mathbb{E}(M)$ , on a
\[ \mathrm{Hom}_k(G,\ \mathbb{E}(M)) = \mathbb{E}(M)(G) = \mathrm{Hom}_{\hand{$\mathcal{D}$}}(M,\ \underset{\rightarrow}{\mathbb{W}}(G))\ , \]
ce qui définira l'isomorphisme désiré si l'on montre qu'un élément
$u \in \mathrm{Hom}_{\hand{$\mathcal{D}$}}(M,\ \underset{\rightarrow}{\mathbb{W}}(G))$ définit un morphisme de groupes $G \to \mathbb{E}(M)$ si et seu-%
lement si u se factorise par Prim(G) .

En effet $\mathrm{Hom}_{k\text{-}Gr}(G,\ \mathbb{E}(M))$ s'identifie à la partie de
$\mathrm{Hom}_{k\text{-}alg}(A,\ \Gamma(\hand{$\underline{\mathcal{O}}$}_G))$ formée des morphismes $\alpha$ tels que le diagramme sui-%
vant soit commutatif :
\diagram{
\begin{tikzcd}
A \arrow[r, "\alpha"] \arrow[d, "\Delta"'] & \Gamma(\hand{$\underline{\mathcal{O}}$}_G) \arrow[d, "\Delta"] \\
A \otimes A \arrow[r, "\alpha \otimes \alpha"] & \Gamma(\hand{$\underline{\mathcal{O}}$}_G) \otimes \Gamma(\hand{$\underline{\mathcal{O}}$}_G)\ .
\end{tikzcd}
}

Mais le morphisme $\Delta : A \to A \otimes A$ est donné par la formule
\[ \Delta(T_x) = S_N(T_{V^N(x)} \otimes 1, \ldots, T_x \otimes 1\ ;\ 1 \otimes T_{V^N(x)}, \ldots,\ 1 \otimes T_x) \]

\margin{53}
